% document formatting
\documentclass[10pt]{article}
\usepackage[utf8]{inputenc}
\usepackage[left=1in,right=1in,top=1in,bottom=1in]{geometry}
\usepackage[T1]{fontenc}
\usepackage{xcolor}
\usepackage[table,xcdraw]{xcolor}

% math symbols, etc.
\usepackage{amsmath, amsfonts, amssymb, amsthm}

% lists
\usepackage{enumerate}
\usepackage{enumitem}
\usepackage{tabularx}
\usepackage{multirow}

% images
\usepackage{graphicx} % for images
\usepackage{tikz}
\usepackage[dvipsnames]{xcolor}

% code blocks
\usepackage{minted, listings} 

% verbatim greek
\usepackage{alphabeta}
\DeclareMathOperator*{\argmin}{arg\,min}
\DeclareMathOperator*{\argmax}{arg\,max}

\newcommand{\dd}{\text{d}}

\graphicspath{{./assets/images/Week 11}}

\title{03-621 Week 11 \\ \large{Advanced Quantative Genetics}}
 
\author{Aidan Jan}

\date{\today}

\begin{document}
\maketitle

\section*{Finding the Genes}
What is the actual genetic architecture of complex traits?
\begin{itemize}
	\item \textbf{Ronald Fisher's Infinitesimal Model (1918)}
	\item Large number of genes, with common alleles each contributing a small additive effect.
\end{itemize}
Does this model hold up in the face of modern GWAS evidence, especially for common threshold disease traits?  How important are rare alleles with large effect?  How important are non-additive effects?

\subsection*{Attempts to Account for a Greater Proportion of Heritability}
We want to account for:
\begin{itemize}
	\item Type II error
	\item Infinitesimal hypothesis
\end{itemize}
Methods:
\begin{itemize}
	\item The basic concept is to combine the effects of \textbf{all} of the SNPs in the study (not just significant SNPs).  Is any $h^2$ really missing?
    \item Several software packages available.  (e.g., GCTA: as random effects by a mixed linear model)
\end{itemize}
The SNP genotypes are used to calculate genetic similarities among unrelated individuals and this is regressed on phenotypic similarities.  Strongly criticized as severely overestimating or underestimateing $h^2$ due to multiple flaws and sensitivities to data imperfections (e.g., population structure due to admixture)

\subsection*{False Positives due to Population Admixture}
When two populations with different trait prevalence intermix, failure to take mixed genetic background into account can lead to false positives in GWAS.
\begin{itemize}
	\item The reason: any allele(s) that distinguish the two populations will appear to be associated with the trait
\end{itemize}

\subsection*{From Mapping to Identifying Genes}
\textbf{General strategies}
\begin{itemize}
	\item Follow-up mapping and association studies at higher resolution
	\item Extend analysis to rare variants in delimited regions
	\item Establish causation
\end{itemize}
\textbf{Some special approaches:} Sometimes useful for monogenic traits and some complex traits
\begin{itemize}
	\item Exome sequencing
	\item Cytological Mapping using chromosome abnormalities
\end{itemize}

\subsection*{Route to Pinpointing the Responsible Gene and Mutations}
\begin{enumerate}
    \item Extend Association Studies to tracing inheritance in families
    \begin{itemize}
        \item Known common and rare alleles from databases
        \item New rare alleles identified by sequencing the candidate genes in patients and controls
    \end{itemize}
    \item Predict effect of mutation on gene function
    \begin{itemize}
        \item From first principles (null alleles, frameshifts, missense mutations)
        \item From knowledge or modeling of structure and functional domains
    \end{itemize}
    \item Verify effect on gene function
    \begin{itemize}
        \item In vitro assays of purified wild type and mutant proteins
        \item In vivo assays in cell lines or model organisms
    \end{itemize}
    \item Verify effect on phenotype
    \begin{itemize}
        \item Gene knockout or replacement in cell lines or model organisms
    \end{itemize}
\end{enumerate}

\subsection*{Cytological Mapping Using Chromosomal Abnormalities}
\begin{center} 
	\includegraphics*[width=0.8\textwidth]{W11_1.png} 
\end{center}
The banding patterns of chromosomes allow us to identify the location of breakpoints in chromosome abnormalities that may \textbf{inactive} or \textbf{inappropriately activate} a gene.  Useful for specific cases of:
\begin{itemize}
	\item Rare disease traits with insufficient number of informative pedigrees
	\item Highly deleterious dominant traits
	\begin{itemize}
        \item Not transmitted
        \item All cases are new sporadic mutations
    \end{itemize}
\end{itemize}

\subsubsection*{Abnormalities 1: Inversions and Deletions}
Genes at the junctions may be \textbf{inactivated} or \textbf{fused}.  Fusing is a common mechnanism of activation of cancer-promoting genes.
\begin{center} 
	\includegraphics*[width=0.75\textwidth]{W11_2.png} 
\end{center}

\subsubsection*{Abnormalities 2: Translocations}
Result from breaks on different chromosomes.
\begin{center} 
	\includegraphics*[width=0.75\textwidth]{W11_3.png} 
\end{center}

\subsection*{Examples of Copy Number Polymorphisms (CNPs) Associated with Complex Diseases}
\begin{center}
    \begin{tabular}{|l|l|}
        \hline
        \textbf{CNP Allele} & \textbf{Disease} \\ \hline
        \begin{tabular}[c]{@{}l@{}}Deletion upstream of \textit{IRGM}\\ (innate immune response)\end{tabular} & Crohn's disease \\ \hline
        Low-copy of repeat in Lipoprotein A gene & Coronary artery disease \\ \hline
        \begin{tabular}[c]{@{}l@{}}Extra copies of the $\beta$-defensin multigene family\\ (antimicrobial peptides)\end{tabular} & Psoriasis \\ \hline
        Deletion of 1 copy of the complement \textit{C4} gene & \multirow{2}{*}{Systemic lupus erythematosus} \\ \cline{1-1}
        Reduced copies of \textit{FCGR3A} (Fc portion of IgG) &  \\ \hline
    \end{tabular}
\end{center}

\subsection*{Examples of Loci that Exhibit Copy Number Variation in Autism Spectrum Disorder}
\begin{center}
    \begin{tabular}{|l|l|l|l|l|}
        \hline
        \textbf{Locus Size} & \textbf{Genes} & \textbf{Location} & \textbf{Pathogenic Allele} & \textbf{Frequency in ASD (\%)} \\ \hline
        0.7 Mb & 30 genes & 16p11.2 & deletions and duplications & 0.8 \\ \hline
        $\sim 1$ Mb & \begin{tabular}[c]{@{}l@{}}\textit{PTCHD1AS}\\ \textit{PTCHD1}\end{tabular} & Xp22.1 & \begin{tabular}[c]{@{}l@{}}deletions affecting \textit{PTCHD1AS}\\ (upstream antisense ncRNA)\end{tabular} & 0.5 \\ \hline
        Variable & \textit{NRXN1} & 2p16.3 & mostly deletions & 0.4 \\ \hline
        1.4 Mb & 22 genes & 7q11.2 & duplication & 0.2 \\ \hline
        2.5 Mb & 56 genes & 22q11.2 & deletions and duplications & 0.2 \\ \hline
        1.5 Mb & 14 genes & 1q21.1 & duplications & 0.2 \\ \hline
        Variable & \textit{shank2} & 11q13.3 & deletions & 0.1 \\ \hline
    \end{tabular}
\end{center}
Note that the copy number variations each account for only a \textbf{small} percentage of ASD cases

\section*{Exome Sequencing}
Also useful for single-gene disorders that are:
\begin{itemize}
	\item Very rare disease traits with an insufficient number of informative pedigrees for mapping
	\item Highly deleterious dominant traits
	\begin{itemize}
        \item not transmitted
        \item therefore, all cases are \textit{new} sporadic mutations
        \item therefore, cannot be mapped
    \end{itemize}
\end{itemize}

\subsection*{Exome sequencing method}
Targets sequencing to exons.
\begin{center} 
	\includegraphics*[width=\textwidth]{W11_4.png} 
\end{center}
Exon-seq reads are aligned to the genome and compared to identify mutations (e.g., rare dominant)
\begin{center} 
	\includegraphics*[width=\textwidth]{W11_5.png} 
\end{center}

\subsection*{Examples of Successful Exome Sequencing}
Very rare recessive or severe dominant monogenic diseases
\begin{center}
    \vspace{-1em}
    \begin{tabular}{|l|c|c|c|}
        \hline
        \multicolumn{1}{|c|}{\textbf{Disease}} & \textbf{Inheritance} & \textbf{Type of Disorder} & \textbf{\begin{tabular}[c]{@{}c@{}}Exomes\\ Sequenced From\end{tabular}} \\ \hline
        Miller Syndrome & AR & \multirow{3}{*}{\begin{tabular}[c]{@{}c@{}}congenital disorders\\ of development\end{tabular}} & \multirow{3}{*}{\begin{tabular}[c]{@{}c@{}}unrelated\\ sporadic cases\end{tabular}} \\ \cline{1-2}
        Kabuki Syndrome & AD &  &  \\ \cline{1-2}
        Schinzel-Giedion Syndrome & AD &  &  \\ \hline
        Osteogenesis imperfecta Type VI & AR & connective tissue disorder & \multirow{2}{*}{\begin{tabular}[c]{@{}c@{}}affected siblings\\ born to consanguineous\\ parents\end{tabular}} \\ \cline{1-3}
        Spastic paraplegia 30 & AD & \begin{tabular}[c]{@{}c@{}}early-onset neuromuscular\\ disorder\end{tabular} &  \\ \hline
    \end{tabular}
\end{center}

\subsubsection*{Exome Sequencing of Siblings to Find the Genetic Basis of the Miller Syndrome}
\begin{enumerate}
    \item Find variants in patients
    \item Filter for homozygous alleles (trait is recessive)
    \item Remove common variants found in many people.
    \item Compare to mutations in other unrelated affected children
\end{enumerate}
\begin{center} 
	\includegraphics*[width=0.7\textwidth]{W11_6.png} 
\end{center}
11 different LOF mutations in the \textit{DHODH} gene found among 6 families with high rate of Miller syntdrome
\begin{itemize}
	\item \textit{DHODH}: Biosynthesis of pyrimidine nucleotides: C, T, U for DNA, RNA
\end{itemize}

\subsection*{Other Uses of Exome-Seq}
Exome-seq is also useful for phenotypes with \textbf{extremely heterogeneous genetics}.  e.g., Autosomal Recessive Intellectual Disabilities (ARID)
\begin{itemize}
	\item Note the involvement in shared functional networks:
	\begin{itemize}
        \item additional proof of relevance
        \item clues to identify additional relevant genes
        \item clues to underlying molecular pathology
    \end{itemize}
\end{itemize}

\subsection*{Limitations of Exome Sequencing}
\begin{itemize}
	\item Relevant mutations are \textbf{not necessarily} in Exons!
	\item Relevant mutations are \textbf{not necessarily} in protein-coding genes at all!
	\item Generally, not useful for complex traits
\end{itemize}

\end{document}